% =====================================================================
% Ch3 report — Recurrent networks on SmartWatch Gestures (1 A4)
% YOU (the student) write ALL prose. This file only provides structure.
%
% WRITE:
%   - Hypothesis: GRU vs LSTM, or Conv1D-before-RNN, grounded in theory
%     (gates, hidden state, 1D convolution on timeseries).
%   - Experiment design: baseline GRU h=64 l=1; what varied (hidden_size,
%     num_layers, dropout, GRU/LSTM, Conv1D+GRU); MLflow tags; random
%     baseline = 5% (1/20 classes).
%   - Results: architecture comparison, accuracy progression / final accuracy,
%     EXPLICITLY showing >90% achieved.
%   - Conclusions.
% REFLECT ON:
%   - Which architecture won and WHY, linked to the data (variable-length,
%     3-axis accelerometer).
%   - Why Conv1D before the RNN makes sense for accelerometer signals.
%   - GRU vs LSTM trade-off (gate count / parameters) on this data size.
%   - What did NOT work (e.g. dropout with num_layers=1 is a no-op) and why.
% =====================================================================
\documentclass[11pt,a4paper]{article}
\usepackage[margin=2cm]{geometry}
\usepackage{graphicx}
\usepackage{float}
\usepackage{caption}
\usepackage{booktabs}
\usepackage{fancyhdr}
\usepackage{parskip}
\setlength{\parskip}{4pt}
\usepackage{titlesec}
\titlespacing*{\section}{0pt}{6pt}{3pt}
\captionsetup{font=small,skip=3pt}

\newcommand{\studentname}{Vincent Brand}
\newcommand{\studentnumber}{1447806}
\newcommand{\course}{Machine Learning -- MADS, HAN}
\newcommand{\reportdate}{2026}
\newcommand{\repourl}{https://github.com/vincentbrand/MADS-ML-Vincent}

\pagestyle{fancy}
\fancyhf{}
\lhead{\studentname\ (\studentnumber)}
\rhead{\course}
\cfoot{\thepage}
\renewcommand{\headrulewidth}{0.4pt}
\sloppy

\newcommand{\figorbox}[2][\linewidth]{%
  \IfFileExists{#2}{\includegraphics[width=#1]{#2}}%
  {\fbox{\parbox[c][3.5cm][c]{0.98\linewidth}{\centering\ttfamily
   missing figure: #2\\[2pt] run the notebook + export step to generate it}}}}

\begin{document}

\begin{center}
  {\Large\bfseries Recurrent networks on SmartWatch gestures}\\[3pt]
  \studentname\ \textbar\ \studentnumber\ \textbar\ \course\\
  \reportdate\ \textbar\ \repourl
\end{center}

\section{Hypothesis}
The gestures are short, variable-length sequences of 3-axis acceleration, and a
gesture is essentially a sequence of \emph{local} motion primitives (a flick, a
turn) rather than one long-range dependency. I therefore hypothesise that (a) a
1D-convolutional front-end before the recurrent layer will help, because Conv1D
extracts these local temporal features that the GRU then orders in time, and (b)
GRU and LSTM will perform almost equally, because the sequences are short enough
that the LSTM's extra cell-state gate buys little over the GRU's lighter
update/reset gates. Random guessing is 5\% (1/20 classes), so any working model
should far exceed that.

\section{Experiment design}
The baseline is the provided GRU (\texttt{hidden\_size}=64, \texttt{num\_layers}=1),
trained with \texttt{CrossEntropyLoss}, Adam and \texttt{ReduceLROnPlateau} for up
to 100 epochs with early stopping (patience 10). From there I varied hidden size
(64/128/256), depth (1/2/3 layers), dropout ($0\to0.3$) and the architecture
family: plain GRU, LSTM, and a Conv1D+GRU hybrid (which \texttt{permute}s
$(B,T,F)\!\to\!(B,F,T)$ for the convolution and back before the GRU). Every run is
logged to MLflow under a descriptive \texttt{model} tag, and Fig.~\ref{fig:arch}
compares their final validation accuracy.

\section{Results}
All six architectures comfortably passed the 90\% target (Fig.~\ref{fig:arch}):
even the GRU baseline reached \textbf{97.97\%}, roughly $20\times$ the 5\% random
floor. The best model is the \textbf{Conv1D+GRU} (conv\_filters=128, kernel=5,
$h$=256, 2 layers) at \textbf{99.53\%} (test loss 0.033), ahead of the best LSTM
(98.91\%) and the best plain GRU (98.59\%). Adding the convolutional front-end
gave the single largest jump over a same-sized GRU, while GRU and LSTM landed
within $\approx0.3$ percentage points of each other. Scaling the GRU from $h$=64,
1 layer to $h$=128, 2 layers helped (97.97\%$\to$98.59\%) but with clearly
diminishing returns.

\begin{figure}[H]
  \centering
  \figorbox[0.66\linewidth]{../figures/ch3-architecture-comparison.png}
  \caption{Final validation accuracy per architecture on the gestures dataset,
  coloured by family (GRU / LSTM / Conv1D+GRU). All models clear the 90\% target;
  the Conv1D+GRU hybrid is best.}
  \label{fig:arch}
\end{figure}

\section{Conclusions}
Both hypotheses held. (a) The Conv1D+GRU won, which fits the data: accelerometer
gestures are built from local motion shapes, and convolution captures those
cheaply before the GRU models their order. (b) GRU $\approx$ LSTM here, so the
LSTM's extra cell state did not pay off on these short sequences and the cheaper
GRU is the better default. The main caveat is that the task is easy --- even the
baseline hit 98\%, so differences between architectures are small; a harder
cross-user split or a tighter epoch budget would separate them better. One thing
that did \emph{not} work as expected: dropout is a no-op with
\texttt{num\_layers}=1, since PyTorch only applies it \emph{between} stacked RNN
layers.

\end{document}
